\section{Data and Splits}
\label{app:dataset}

\subsection{Filtering and Split Construction}

% 原始数据包含八个域内任务维度；经过任务筛选和样本质量过滤后保留七个，
% 排除 \texttt{verifiability\_extraction}。四个 RevUtil 任务分别从
% 4,800 条候选样本开始，并删除 human-hard、与测试集重叠以及包含异常
% 字符的记录。两个数据源的训练池均通过教师蒸馏构建，仅保留格式有效且
% 教师分数与 gold label 一致的轨迹。测试集清洗还删除了两条重复的
% Coherence 样本，以及两条与训练集重叠的 Positioning Check 样本。
% 教师蒸馏的具体设置见附录~\ref{app:teacher-distillation}。
The source contains eight in-domain task dimensions; task and sample-quality
filtering retained seven, excluding \texttt{verifiability\_extraction}.
For each of the four RevUtil tasks, a 4,800-example candidate pool was
constructed after removing human-hard examples, test-set overlaps, and
malformed-character records. Both data sources use training pools constructed
through teacher distillation, retaining only trajectories with valid formatting
and teacher scores matching the gold labels. Test-set cleanup additionally
removed two duplicate Coherence examples and two training-overlap Positioning
Check examples. Teacher-distillation details are reported in
Appendix~\ref{app:teacher-distillation}.

% 所有 prompt 均删除完整的 \texttt{[EXAMPLES]} 区块，以及 system prompt
% 中对应的示例说明句；原始 query、criteria、评分规则和待评价文本保持
% 不变。同时为样本添加稳定标识符和监督元数据。每个样本都被构造成严格
% 对齐的理由与分数监督（RS）和仅分数监督（SO）格式。两种格式共享样本、
% 输入、gold label 和数据划分，仅 system instruction 与 completion 不同。
Every prompt was stripped of the complete \texttt{[EXAMPLES]} block and the
corresponding system-prompt sentence, while the original query, criteria,
scoring rubric, and evaluated instance were preserved. Stable identifiers and
supervision metadata were added. Each example was instantiated in strictly
aligned rationale-and-score (RS) and score-only (SO) formats. The formats share
examples, inputs, gold labels, and data splits; only the system instructions
and completions differ.

% 表~\ref{tab:data-filter-coverage} 展示各任务如何从原始数据经过候选样本
% 构建和质量筛选，形成本文最终使用的数据集。表~
% \ref{tab:data-label-distributions} 进一步给出所得数据在训练池和测试集中的
% 具体标签分布。
Table~\ref{tab:data-filter-coverage} shows how the source data for each task
are converted through candidate construction and quality filtering into the
final dataset used in this study. Table~\ref{tab:data-label-distributions}
further reports the label distributions of the resulting training pools and
test sets.

\begin{table}[H]
\centering
\fontsize{9.5pt}{10.8pt}\selectfont
\renewcommand{\arraystretch}{0.88}
\setlength{\tabcolsep}{1.8pt}
\begin{tabular*}{\columnwidth}
  {@{\extracolsep{\fill}}lrrr@{}}
\toprule
Task & Source & Cand./attempt. & Final \\
\midrule
Coherence & 4,890 & 1,938 & 1,526 \\
Positioning Check & 2,822 & 2,822 & 2,666 \\
Positioning Type & 954 & 954 & 944 \\
Actionability & 10,432 & 4,800 & 1,788 \\
Grounding Spec. & 10,431 & 4,800 & 2,652 \\
Helpfulness & 10,430 & 4,800 & 2,279 \\
Verifiability & 8,323 & 4,800 & 1,825 \\
\bottomrule
\end{tabular*}

% 各任务的源数据量、候选或实际处理数量及最终保留数量。
\caption{Source, candidate/attempted, and retained counts by task.}
\label{tab:data-filter-coverage}
\end{table}

\begin{table}[H]
\centering
\fontsize{9.5pt}{10.8pt}\selectfont
\renewcommand{\arraystretch}{0.92}
\setlength{\tabcolsep}{1.5pt}

\begin{tabular*}{\columnwidth}
  {@{\extracolsep{\fill}}lrrrr@{}}
\toprule
& \multicolumn{2}{c}{Train pool}
& \multicolumn{2}{c}{Test} \\
\cmidrule(lr){2-3}\cmidrule(lr){4-5}
Binary task & Label 0 & Label 1 & Label 0 & Label 1 \\
\midrule
Coherence & 770 & 756 & 523 & 523 \\
Positioning Check & 1,499 & 1,167 & 341 & 262 \\
Positioning Type & 628 & 316 & 136 & 68 \\
\bottomrule
\end{tabular*}

\vspace{3pt}
\begin{tabular*}{\columnwidth}
  {@{\extracolsep{\fill}}lrrrrr@{}}
\toprule
\multicolumn{6}{c}{\textbf{Ordinal train-pool labels}} \\
\midrule
Task & 1 & 2 & 3 & 4 & 5 \\
\midrule
Actionability & 148 & 361 & 521 & 86 & 672 \\
Grounding Spec. & 104 & 78 & 839 & 66 & 1,565 \\
Helpfulness & 122 & 358 & 853 & 858 & 88 \\
Verifiability & 964 & 247 & 373 & 172 & 69 \\
\bottomrule
\end{tabular*}

\vspace{3pt}
\begin{tabular*}{\columnwidth}
  {@{\extracolsep{\fill}}lrrrrr@{}}
\toprule
\multicolumn{6}{c}{\textbf{Ordinal test labels}} \\
\midrule
Task & 1 & 2 & 3 & 4 & 5 \\
\midrule
Actionability & 234 & 164 & 312 & 95 & 195 \\
Grounding Spec. & 123 & 64 & 389 & 41 & 383 \\
Helpfulness & 34 & 107 & 462 & 311 & 86 \\
Verifiability & 221 & 77 & 346 & 120 & 24 \\
\bottomrule
\end{tabular*}

% 最终训练池和测试集中的标签分布。
\caption{Label distributions in the final training pools and test sets.}
\label{tab:data-label-distributions}
\end{table}

\subsection{Regenerated-Rationale Data}
\label{app:reason-training-data}

% 对于两种方法，第二阶段训练数据均由各自方法的第一阶段模型重新生成；
% 第一阶段模型使用其自身的训练目标和 seed 42 进行训练。两种方法采用
% 相同的推理配置，并将生成分数替换为 gold score。七个任务包括三个
% Related Work 数据集和四个 RevUtil 数据集，共包含 13,680 个训练池
% 样本。这些样本使用 seed 20260720 划分为 12,312 个训练样本和
% 1,368 个验证样本。
For both methods, the second-stage training data are regenerated by the
first-stage model of each respective method (trained with its own training
objective, seed 42) under identical inference configurations, with generated
scores replaced by gold scores. The seven tasks cover three Related Work and
four RevUtil datasets, totaling 13,680 training-pool examples, split with seed
20260720 into 12,312 training and 1,368 validation instances.

\subsection{Paired DS/RF Views}
\label{tab:constructing-pair-data}

% PaIR 覆盖七个任务和 13,680 个源样本。每个源样本分别构造一个 DS
% 视图和一个 RF 视图，并将两个视图放入同一个 micro-batch。训练时，
% 通过各视图独立的 loss mask 分别计算两个视图上的 token-level
% cross-entropy；prompt token 的标签被设为 \(-100\)，因此只有
% completion token 参与监督。每个视图的损失分别在其自身的监督 token
% 上取平均，总损失为
% \(0.5L_{\mathrm{DS}} + 0.5L_{\mathrm{RF}}\)。
% 使用 seed 20260720 和 0.1 的验证集比例后，配对数据被划分为
% 12,312 个用于优化训练的样本和 1,368 个验证样本；测试集不参与训练，
% 以保证匹配评估。
PaIR covers seven tasks and 13,680 source examples, each paired with a DS view
and an RF view that appear in the same micro-batch. During training,
token-level cross-entropy is computed over the two views separately via
per-view loss masks (prompt tokens are masked to \(-100\), so only completion
tokens contribute to supervision); each view's loss is averaged over its own
supervised tokens, and the total loss is
\(0.5L_{\mathrm{DS}} + 0.5L_{\mathrm{RF}}\). Using seed 20260720 and a
validation ratio of 0.1, the paired data yield 12,312 optimizer-training and
1,368 validation examples; test sets remain outside training for matched
evaluation.